\Needspace{6\baselineskip}
\begin{center}\hypertarget{native-260-article187}{}\pdfbookmark[2]{187.}{article-187}187.\end{center}
Wir kehren von diesen allgemeinen Untersuchungen zu der besonderen Aufgabe
zurück, derentwegen wir sie angestellt haben. Bevor wir eine Bahn aus mehr
Beobachtungen als unbedingt nötig möglichst genau bestimmen können, muss
bereits eine Näherungsbestimmung vorliegen, die von allen gegebenen
Beobachtungen nur wenig abweicht. Gesucht sind die Korrekturen, die an
diesen Näherungselementen noch anzubringen sind, um eine möglichst genaue
Übereinstimmung zu erzielen. Da man annehmen darf, dass diese Korrekturen
so \SourceLeaf{261}klein ausfallen, dass ihre Quadrate und Produkte vernachlässigt
werden können, lassen sich die daraus folgenden Änderungen der berechneten
geozentrischen Örter des Himmelskörpers mit den Differentialformeln aus
Abschnitt II des ersten Buches berechnen. Die mit den gesuchten korrigierten
Elementen berechneten Örter werden somit durch lineare Funktionen der
Korrekturen an den Elementen dargestellt. Der Vergleich mit den beobachteten
Örtern nach den oben dargelegten Grundsätzen führt zur Bestimmung der
wahrscheinlichsten Werte. Diese Rechenschritte sind so einfach, dass sie
keiner weiteren Erläuterung bedürfen. Es liegt auf der Hand, dass man
beliebig viele Beobachtungen heranziehen kann, wie weit sie auch
auseinanderliegen. --- Dasselbe Verfahren lässt sich auch zur Korrektur
\emph{parabolischer} Kometenbahnen verwenden, wenn eine längere
Beobachtungsreihe vorliegt und eine möglichst gute Übereinstimmung
verlangt wird.
